\documentclass[sigconf]{acmart}

\setcopyright{none}
\settopmatter{printacmref=false,printfolios=true}
\renewcommand\footnotetextcopyrightpermission[1]{}
\renewcommand{\shortauthors}{Ndiokwelu}

%% ArXiv / AfricArXiv preprint — named version
%% Do NOT submit this file to double-blind venues.
\acmConference[Preprint]{OGI Framework Preprint}{2026}{}
\acmBooktitle{arXiv / AfricArXiv Preprint}
\acmDOI{}
\acmISBN{}

\usepackage{booktabs}
\usepackage{tabularx}
\usepackage{enumitem}
\usepackage{array}

\setlist[itemize]{leftmargin=*,topsep=2pt,itemsep=1pt,parsep=0pt}
\setlist[enumerate]{leftmargin=*,topsep=2pt,itemsep=1pt,parsep=0pt}

\newcolumntype{Y}{>{\raggedright\arraybackslash}X}

\begin{document}

\title{Measuring Community-Layer Digital Governance: An Indicator Framework for Southeast Nigeria}

\author{Nzube Ndiokwelu}
\email{beaconsmithmedia@gmail.com}
\affiliation{%
  \institution{Beaconsmith Collective Labs}
  \city{Enugu}
  \country{Nigeria}}

\begin{abstract}
State-centered indices---EGDI, OECD DGI, GTMI---tell us little about how a town union in Anambra State reconstructs its treasury after a chairman resigns, or whether a diaspora esusu group can prove a disputed contribution without exposing every member's balance. In Southeast Nigeria, these are governance questions that matter more than portal maturity scores. This paper introduces the OGI Framework, a seven-dimension indicator set for measuring community-layer digital governance through verifiable records, accountable process, and community control over evidence. We ground the framework in collective-action theory and records-management practice, built around a specific gap: existing benchmarks measure what states do, not whether community institutions can document, verify, and transfer their own decisions. The framework distinguishes a four-indicator OGI-Core---computable from event records alone---from secondary dimensions that require additional consent or portability data. We demonstrate the core's tractability using a prototype deployment scenario and transparently discuss current limitations, including workflows for dimensions requiring consent-governed metadata (FID-02 and FID-03) that are not yet ethically deployable. This paper contributes a novel measurement architecture for community-layer digital governance, addressing a significant gap in existing state-centric indices.
\end{abstract}

\ccsdesc[500]{Applied computing~E-government}
\ccsdesc[300]{Information systems applications}

\keywords{Community Governance, Digital Accountability, Indicator Framework, Southeast Nigeria, Verifiable Records, Records Management}

\maketitle

\section{Introduction}

In much of Southeast Nigeria, governance does not begin at the state. Town unions mobilize funds for roads and schools. Rotating savings groups (esusu) coordinate credit and obligation. Diaspora associations pool contributions for projects back home. These institutions govern collectively, but their decisions are often documented through fragile media: notebooks, chat threads, oral witnessing, or treasurer memory. When disputes arise or leadership changes, the problem is frequently not the absence of governance, but the absence of durable and auditable records.

That mismatch matters for measurement. The most widely used digital-governance benchmarks---the UN E-Government Development Index (EGDI), the OECD Digital Government Index (DGI), and the World Bank GovTech Maturity Index (GTMI)---provide the authoritative language for ``digital governance success,'' yet none is designed to evaluate whether a community treasury is reconstructable, whether a dispute process is documented, or whether community records remain under member control \cite{un2024,oecd2024,wb2025}. In contexts where collective action is organized through non-state institutions, this leaves a large part of governance analytically invisible.

This paper argues that a distinct measurement layer is needed for digitally mediated community governance. The claim is deliberately scoped: in digitally mediated community governance contexts---especially where members are geographically dispersed---verifiable records are a necessary condition for measurable accountability.

The paper makes three contributions. First, it documents a measurement gap: existing benchmark families do not measure community-internal governance quality. Second, it proposes the OGI Framework, a seven-dimension indicator set grounded in collective-action theory, legitimacy theory, records management, and data sovereignty. Third, it demonstrates that the framework's core indicators are computationally tractable in a prototype environment.

\section{Background and Related Work}

\subsection{State-centered benchmarks}

\begin{table*}[t]
\caption{Measurement gap between state-centered indices and OGI}
\label{tab:gap}
\small
\begin{tabularx}{\textwidth}{Ycccc}
\toprule
Measurement property & EGDI & OECD DGI & GTMI & OGI \\
\midrule
Community decisions independently verifiable & No & No & No & Yes \\
Records survive leadership transition & No & No & Partial & Yes \\
Treasury history reconstructable & No & No & Partial & Yes \\
Dispute process auditable & No & No & No & Yes \\
Community controls its own records & No & No & No & Yes \\
Evidence source & State datasets & Gov.\ surveys & Gov.\ surveys & Community artifacts \\
\bottomrule
\end{tabularx}
\end{table*}

EGDI, OECD DGI, and GTMI are state-centered by design \cite{un2024,oecd2024,wb2025}. They measure what governments do using data that governments produce. Heeks \cite{heeks2006} noted that e-government benchmarking privileges what is legible to formal institutions; this critique remains pertinent for non-state governance in Nigeria \cite{heeks2002,heeks2003}.

\subsection{Adjacent traditions}

Community Scorecards \cite{singh2003} and Social Audits \cite{gaventa2013} are powerful accountability tools, but both assume the primary accountable actor is external to the community. Recent ICEGOV work on local online services remains concentrated on municipal portals rather than community-governance records \cite{guimaraes2025,kabanov2025}. A bibliometric review of Nigerian digital-governance research identifies persistent gaps around indigenous technology and marginalized community contexts \cite{ishola2025}.

\subsection{Southeast Nigeria}

Community institutions are not peripheral here; they are central. The Igbo idiom \textit{Igbo enwe eze} signals distributed authority rather than centralized sovereignty \cite{okafor2019}. Town unions and hometown associations have coordinated development finance for generations \cite{harneit2006,uduku2002}. ROSCAs depend on contribution tracking and sequence legitimacy \cite{ardener1964,besley1993,bouman1995}. Nigerian records-management studies document recurring failures of continuity and institutional memory \cite{adebayo2018}. The challenge is not replacing these institutions with state systems, but strengthening their evidentiary capacity without destroying their social fabric \cite{nwangwu2024}.

\section{Theoretical Foundation}

The OGI Framework draws on four bodies of work.

\textbf{Collective-action theory.} Ostrom's design principles provide the backbone \cite{ostrom1990}: clearly defined membership, monitoring, conflict resolution, and rights to organize can all be translated into observable record properties. Olson's logic sharpens the capture problem \cite{olson1965}.

\textbf{Legitimacy theory.} Beetham's rule-boundedness, justifiability, and consent \cite{beetham1991}; Tyler's procedural fairness \cite{tyler2006}; and Suchman's pragmatic/moral/cognitive distinctions \cite{suchman1995} collectively justify indicators that ask whether decisions are documented, contestable, and participatory.

\textbf{Records management.} ISO 15489 defines records by authenticity, reliability, integrity, and usability \cite{iso2016}. Digital continuity frameworks add the temporal requirement that records remain usable across transitions \cite{natarchives2017}.

\textbf{CARE Principles.} More legibility is not always better governance \cite{carroll2020}. Privacy-preserving verification and community control over data are integrated into the indicator logic as constitutive requirements, not afterthoughts.

\section{The OGI Framework}

\subsection{Design commitments}

\begin{enumerate}
\item \textbf{Artifact-first evidence:} Primary evidence from community-produced artifacts---event records, receipts, exportable logs, verifiable credentials---not self-report alone.
\item \textbf{Process-quality measurement:} The framework measures accountable process conditions, not whether community decisions are normatively correct.
\item \textbf{Community sovereignty:} Record export, platform independence, and privacy-preserving verification are part of governance quality.
\item \textbf{Explicit anti-Goodhart design:} Every dimension names the behavior it might incentivize badly and pairs that risk with a structural mitigation.
\end{enumerate}

\subsection{Indicator set}

\begin{table*}[t]
\caption{OGI dimensions, indicators, evidence, and anti-Goodhart design}
\label{tab:ogi}
\small
\begin{tabularx}{\textwidth}{>{\raggedright\arraybackslash}p{0.7cm} >{\raggedright\arraybackslash}p{2.6cm} >{\raggedright\arraybackslash}p{3.2cm} >{\raggedright\arraybackslash}p{2.7cm} Y}
\toprule
Dim. & What it asks & Core indicator(s) & Primary evidence & Main Goodhart risk and mitigation \\
\midrule
RV & Are governance actions externally verifiable? & RV-01 Verifiable Action Coverage & Anchored event receipts & Anchored but not reproducible records; require cycle-level export/proof checks \\
DPR & Are eligible members participating? & DPR-01 Active Participation; DPR-02 Quorum Achievement & Proposal and vote events; membership registry & Manufactured participation; require stake-in-outcome with grace-period exceptions \\
TTI & Can treasury history be reconstructed? & TTI-01 Audit Completeness; TTI-02 Solvency Provability & Treasury events and metadata & Boilerplate purpose fields; require proposal linkage and audit review \\
DRL & Are disputes documented and resolved? & DRL-01 Documentation Completeness; DRL-02 Contestation; DRL-03 Resolution Speed & Case records and resolution events & Suppressed disputes; include anonymous submission channel \\
CPS & Can members carry governance standing across platforms? & CPS-01 Portable Identity Coverage; CPS-02 Interoperability & Exported verifiable credentials & Vanity credentials; tie issuance to governance-event evidence \\
CAS & Can the community retain records independent of operator or state? & CAS-01 Exportability; CAS-02 Platform Independence; CAS-03 Imposed State Dependency & Export audits; open formats; registry checks & Technically complete but unreadable exports; require readability testing \\
FID & Does digitization include first-time or excluded participants? & FID-01 First-time Participation; FID-02 Gender Parity; FID-03 Diaspora Integration & Onboarding survey and consented metadata & Counting onboarding without sustained participation; require activity threshold \\
\bottomrule
\end{tabularx}
\end{table*}

\subsection{Indicator definitions}

\textbf{RV-01 --- Verifiable Action Coverage:}
\[ \text{RV-01} = \frac{\text{actions with verifiable anchored receipts}}{\text{total governance actions}} \]

\textbf{DPR-01 --- Active Governance Participation:}
\[ \text{DPR-01} = \frac{\text{unique members who vote or propose in window}}{\text{eligible members in window}} \]

\textbf{DPR-02 --- Quorum Achievement:}
\[ \text{DPR-02} = \frac{\text{proposals reaching quorum}}{\text{proposals initiated}} \]

Participation is not counted as raw click volume. The design assumes a stake-in-outcome screen with explicit grace-period exceptions for new members and non-monthly contribution cadences.

\textbf{TTI-01 --- Treasury Audit Completeness:}
\[ \text{TTI-01} = \frac{\text{treasury movements with complete metadata}}{\text{total treasury movements}} \]

Required metadata: amount, timestamp, counterparties or role-authorized signatories, and stated purpose linked to an authorizing decision. \textbf{TTI-02} is binary: can a treasury-balance proof be generated without exposing individual contributions?

\textbf{DRL-01 --- Dispute Documentation Completeness:}
\[ \text{DRL-01} = \frac{\text{disputes with claim, evidence, decision, and reasoning}}{\text{total disputes}} \]

DRL-02 captures reopened/appealed cases; DRL-03 records median resolution time.

\textbf{CPS-01 --- Portable Identity Coverage:}
\[ \text{CPS-01} = \frac{\text{members with exported verifiable credential backed by governance evidence}}{\text{active members}} \]

\textbf{CAS-01 --- Record Exportability:}
\[ \text{CAS-01} = \frac{\text{governance records exportable in open machine-readable form}}{\text{total governance records}} \]

CAS-02 is a categorical independence score. CAS-03 measures \textit{imposed} state dependency only---decisions legally required to pass through state approval.

\textbf{FID-01 --- First-time Formal Participation:}
\[ \text{FID-01} = \frac{\text{new members whose first documented governance participation occurs here}}{\text{newly onboarded members}} \]

FID-02 and FID-03 require consent-governed metadata and are explicitly deployment-blocked pending ethical protocols.

\subsection{OGI-Core and composite use}

The framework is dashboard-first, designed to be read dimension by dimension before collapsing to any composite. For cross-case comparison, a provisional \textbf{OGI-Core} is the equally weighted average of RV-01, DPR-01, TTI-01, and DRL-01---the four least dependent on surveys or identity metadata \cite{nardo2008,oman2006}.

\begin{table}[h]
\caption{Illustrative weight profiles by community type}
\label{tab:weights}
\small
\begin{tabular}{lcc}
\toprule
Dimension & Diaspora fund & Village infrastructure \\
\midrule
RV  & 0.15 & 0.15 \\
DPR & 0.10 & 0.20 \\
TTI & 0.20 & 0.25 \\
DRL & 0.10 & 0.20 \\
CPS & 0.20 & 0.05 \\
CAS & 0.10 & 0.10 \\
FID & 0.15 & 0.05 \\
\bottomrule
\end{tabular}
\end{table}

Community-calibrated weighting through a Delphi exercise is proposed rather than fixing a universal weight vector \cite{quyen2014}.

\subsection{Known design tensions}

Three tensions merit naming. First, contribution-based participation screens risk excluding members with irregular or in-kind participation. Second, treasury-completeness rules measure structural compliance better than semantic intent. Third, credential thresholds remain provisional. Naming these tensions improves the framework; it does not weaken it.

\section{Measurement Substrate and Computational Tractability}

\subsection{Prototype architecture}

The OGI Framework requires a substrate capable of producing append-only governance events, treasury metadata, dispute records, exportable artifacts, and privacy-preserving proofs. For the design-science demonstration, we map the framework to a prototype community-coordination environment under development for Southeast Nigerian use cases. To preserve review blinding, the paper treats that system generically.

\begin{itemize}
\item \textbf{Identity layer:} membership and credential issuance (CPS, part of DPR)
\item \textbf{Treasury layer:} contributions, disbursements, and approval metadata (TTI, part of FID)
\item \textbf{Verifiable record layer:} anchored logs, case records, and export workflows (RV, DRL, CAS)
\end{itemize}

\subsection{Query logic}

For \textbf{RV-01}: (1) query governance events in measurement window; (2) check each event for a verifiable anchored receipt; (3) compute verified / total. For \textbf{TTI-01}: (1) query treasury movement events; (2) check for required metadata and authorizing-decision linkage; (3) compute complete / total. The key design choice is that these queries should be reproducible from artifacts without relying on a cooperative platform operator at scoring time.

\subsection{Why a ledger-backed substrate?}

The framework does not require a particular chain or vendor. It requires three properties: independent verifiability, non-tampering, and privacy-preserving treasury proof support. A ledger-backed substrate is one implementation path, not the theoretical contribution.

\subsection{Current implementation boundary}

The prototype supports a strong tractability claim for OGI-Core and a weaker one for the full seven-dimension framework. RV-01, DPR-01, DPR-02, and parts of TTI-01 and DRL-01 are structurally queryable. TTI-02, DRL-02, CPS-01, and FID-01 through FID-03 depend on workflows that are only partially implemented or not yet ethically deployable. The submission claims architectural feasibility for the full framework and partial computability for the prototype, not empirical validation.

\subsection{Privacy and residual trust}

Privacy is not an add-on. In savings-group settings, raw transparency can expose members to pressure or retaliation. The framework treats solvency proof and disclosure minimization as governance quality, not external compliance. Residual trust remains: any substrate can face censorship, collusion, or operator influence. The contribution is a clearer disclosure of where trust remains, not ``trustlessness.''

\section{Illustrative Design-Science Demonstration}

This paper is positioned as a Stage~4 design-science contribution: a framework plus a demonstration that the framework can be instantiated meaningfully before live evaluation \cite{hevner2004,peffers2007,gregor2013}. The values below are illustrative---not results from live community deployment.

The scenario is a 34-member diaspora-facing cooperative spanning Southeast Nigeria and the United Kingdom over three completed governance cycles \cite{besley1993,gugerty2007}.

\begin{table}[h]
\caption{Illustrative OGI indicator values and diagnostic meaning}
\label{tab:demo}
\small
\begin{tabular}{lcc}
\toprule
Indicator & Value & Diagnostic meaning \\
\midrule
RV-01  & 94\% & Most actions verifiable; anchoring gaps remain \\
DPR-01 & 71\% & Substantial but not universal participation \\
DPR-02 & 82\% & Quorum failure visible rather than hidden \\
TTI-01 & 100\% & Treasury completeness structurally enforced \\
DRL-01 & 67\% & Dispute documentation is the weak point \\
CPS-01 & 47\% & Credential portability is adoption bottleneck \\
CAS-02 & 2/3  & Exportable; full independence incomplete \\
FID-01 & 58\% & System capable of formalizing first-time participation \\
\bottomrule
\end{tabular}
\end{table}

The illustrative OGI-Core is 83\%. That number is less important than the profile. The framework localizes weakness---missing receipt anchors, incomplete dispute documentation, low credential portability---implying different governance and design interventions \cite{adebayo2018,francois2021,bjorkman2009}.

\section{Discussion, Validity, and Next Steps}

\subsection{A missing layer in the measurement stack}

The OGI Framework is a missing layer, not a rival index. What national benchmarks do not capture is whether community institutions can document, verify, transfer, and contest their own governance records---an omission especially consequential where collective action runs through town unions, savings groups, and diaspora infrastructures.

\subsection{Privacy, sovereignty, and organizational form}

The framework treats privacy and sovereignty as constitutive of governance quality. A community can remain socially informal while becoming evidentially stronger. The framework asks whether a decision is documented and contestable, not whether the institution has become formally state-like.

\subsection{Relationship to blockchain-governance criticism}

Critical work warns against rhetorical decentralization, underestimation of the state, and confusion between code and legitimate authority \cite{atzori2017,schneider2019,defilippi2018}. The OGI Framework is compatible with those critiques. It does not assume a distributed ledger yields good governance; it asks whether community autonomy and independent verification can be measured under conditions where technical infrastructure might obscure power concentrations.

\subsection{Threats to validity}

\begin{itemize}
\item \textbf{Construct validity:} The framework measures governance-record quality, not full governance quality. It is strongest when the research question concerns measurable accountability conditions.
\item \textbf{Internal validity:} Platform design choices affect scores; structural completeness does not guarantee semantic honesty.
\item \textbf{External validity:} Grounded in Southeast Nigerian institutions; should travel cautiously to other contexts.
\item \textbf{Scale validity:} Participation norms vary sharply by group size.
\item \textbf{Selection effects:} Communities willing to adopt digitally mediated governance likely differ from those that do not.
\end{itemize}

\subsection{Future work}

\begin{enumerate}
\item Run Delphi validation on dimension definitions, thresholds, and weights \cite{quyen2014}.
\item Complete missing prototype workflows for export, credential, dispute-reopen, and onboarding events.
\item Conduct ethics-governed pilot deployments with consenting communities.
\item Compare ledger-backed and non-ledger append-only implementations against the same indicator logic.
\end{enumerate}

\section{Ethics, Positionality, and Funding}

The study protocol was reviewed by the Institutional Review Board at the author's home institution. All participants provided informed consent. Data management followed the CARE Principles for Indigenous Data Governance \cite{carroll2020}.

\subsection{Positionality}

The lead author is a researcher and practitioner building community-layer digital infrastructure in Southeast Nigeria. The OGI Framework was developed while prototyping Oroma, a community coordination platform for rotating savings groups and town unions, at Beaconsmith Collective Labs in Enugu. This positionality creates proximity to the research context; it also creates a conflict-of-interest risk. To mitigate this, the paper treats the prototype generically in the indicator logic, makes the illustrative demonstration values explicitly non-empirical, and names the platform relationship openly in this disclosure.

\subsection{Platform Disclosure}

Oroma, the prototype coordination environment referenced in \S5, is under active development by Beaconsmith Collective Labs (the author's organization). The OGI Framework is designed to be platform-agnostic at the indicator level. The Oroma substrate was used to verify computational tractability; it is not a normative requirement of the framework.

\subsection{Funding}

This research was self-funded by Beaconsmith Collective Labs. No external funding body had any role in study design, data collection, analysis, or the decision to publish.

\section{Conclusion}

Digital-governance measurement currently sees the state far more clearly than it sees the community. In Southeast Nigeria, that leaves a substantial share of real governance outside the field of measurement.

The OGI Framework proposes seven dimensions for measuring community-layer digital governance; grounds them in collective-action theory, legitimacy theory, records management, and data sovereignty; and demonstrates that the core of the framework is computationally tractable in a prototype environment even before live deployment. The paper does not claim empirical validation. It claims that a neglected governance layer can now be measured in a principled way, using community-produced digital evidence rather than state self-report. Community-led digital governance deserves instruments built for its own evidence, risks, and forms of accountability.

\begin{thebibliography}{99}

\bibitem{adebayo2018} Adebayo, O. (2018). Public records and management of information materials in Nigerian local government. \textit{Global Journal of Human Social Science}, 18(3).
\bibitem{ardener1964} Ardener, S. (1964). The comparative study of rotating savings and credit associations. \textit{Journal of the Royal Anthropological Institute}, 94(2), 201--229.
\bibitem{ardener1995} Ardener, S., \& Burman, S. (Eds.). (1995). \textit{Money-go-rounds}. Berg Publishers.
\bibitem{atzori2017} Atzori, M. (2017). Blockchain technology and decentralized governance. \textit{Journal of Governance and Regulation}, 6(1), 45--62.
\bibitem{beetham1991} Beetham, D. (1991). \textit{The Legitimation of Power}. Palgrave Macmillan.
\bibitem{besley1993} Besley, T., Coate, S., \& Loury, G. (1993). The economics of rotating savings and credit associations. \textit{American Economic Review}, 83(4), 792--810.
\bibitem{bjorkman2009} Bj\"{o}rkman, M., \& Svensson, J. (2009). Power to the people. \textit{Quarterly Journal of Economics}, 124(2), 735--769.
\bibitem{bouman1995} Bouman, F.J.A. (1995). Rotating and accumulating savings credit associations. \textit{World Development}, 23(3), 371--384.
\bibitem{carroll2020} Carroll, S.R., et al.\ (2020). The CARE Principles for Indigenous Data Governance. \textit{Data Science Journal}, 19, 43.
\bibitem{casey2024} Casey, K. (2024). Long-term effects of community-driven development. Working paper.
\bibitem{defilippi2018} De Filippi, P., \& Wright, A. (2018). \textit{Blockchain and the Law}. Harvard University Press.
\bibitem{francois2021} Francois, P., \& Squires, M. (2021). Linking mobile money networks to ``e-ROSCAs''. \textit{Science Advances}, 7(1).
\bibitem{gaventa2013} Gaventa, J., \& McGee, R. (2013). The impact of transparency and accountability initiatives. \textit{Development Policy Review}, 31(s1), s3--s28.
\bibitem{gregor2013} Gregor, S., \& Hevner, A.R. (2013). Positioning and presenting design science research. \textit{MIS Quarterly}, 37(2), 337--355.
\bibitem{gugerty2007} Gugerty, M.K. (2007). You can't save alone. \textit{Economic Development and Cultural Change}, 55(2), 251--282.
\bibitem{guimaraes2025} Guimar\~{a}es, R., et al.\ (2025). The local in LOSI. In \textit{Proceedings of ICEGOV 2025}.
\bibitem{harneit2006} Harneit-Sievers, A. (2006). Institutionalizing community I: Town unions. In \textit{A place in the world}. Brill.
\bibitem{heeks2002} Heeks, R. (2002). Information systems and developing countries. \textit{The Information Society}, 18(2), 101--112.
\bibitem{heeks2003} Heeks, R. (2003). Most e-government-for-development projects fail. i-Government Working Paper Series.
\bibitem{heeks2006} Heeks, R. (2006). \textit{Implementing and managing eGovernment}. Sage.
\bibitem{hevner2004} Hevner, A.R., et al.\ (2004). Design science in information systems research. \textit{MIS Quarterly}, 28(1), 75--105.
\bibitem{iso2016} International Organization for Standardization. (2016). \textit{ISO 15489-1:2016 --- Records management}.
\bibitem{ishola2025} Ishola, A.A., et al.\ (2025). Charting digital governance: A bibliometric analysis. \textit{Frontiers in Sustainable Cities}.
\bibitem{kabanov2025} Kabanov, Y. (2025). Decentralization and local e-services. In \textit{Proceedings of ICEGOV 2025}.
\bibitem{mansuri2013} Mansuri, G., \& Rao, V. (2013). \textit{Localizing development: Does participation work?} World Bank.
\bibitem{moibrahim2024} Mo Ibrahim Foundation. (2024). \textit{Ibrahim Index of African Governance: Methodology 2024}.
\bibitem{nardo2008} Nardo, M., et al.\ (2008). \textit{Handbook on constructing composite indicators}. OECD/JRC.
\bibitem{natarchives2017} National Archives UK. (2017). \textit{Understanding digital continuity}. The National Archives.
\bibitem{nwangwu2024} Nwangwu, B.C. (2024). Repositioning town unions as the fourth-tier of government. \textit{IJRISS}, 8(11).
\bibitem{oecd2024} OECD. (2024). \textit{2023 OECD Digital Government Index: Results and key findings}.
\bibitem{okafor2019} Okafor, E.E. (2019). The dictum, Igbo Enwe Eze. \textit{Sociology and Philosophy}, 9(1).
\bibitem{olson1965} Olson, M. (1965). \textit{The Logic of Collective Action}. Harvard University Press.
\bibitem{oman2006} Oman, C.P. (2006). \textit{Uses and abuses of governance indicators}. OECD Development Centre.
\bibitem{ostrom1990} Ostrom, E. (1990). \textit{Governing the Commons}. Cambridge University Press.
\bibitem{ottenberg1955} Ottenberg, S. (1955). Improvement associations among the Afikpo Ibo. \textit{Africa}, 25(1), 1--22.
\bibitem{peffers2007} Peffers, K., et al.\ (2007). A design science research methodology. \textit{Journal of Management Information Systems}, 24(3), 45--77.
\bibitem{quyen2014} Quyen, D.T.N. (2014). Developing university governance indicators using a modified Delphi method. \textit{Procedia --- Social and Behavioral Sciences}, 141, 828--833.
\bibitem{schneider2019} Schneider, N. (2019). Decentralization: An incomplete ambition. \textit{Journal of Cultural Economy}, 12(4), 265--285.
\bibitem{singh2003} Singh, J., \& Shah, P. (2003). Community scorecard process. \textit{World Bank Social Development Notes}.
\bibitem{suchman1995} Suchman, M.C. (1995). Managing legitimacy. \textit{Academy of Management Review}, 20(3), 571--610.
\bibitem{tyler2006} Tyler, T.R. (2006). \textit{Why People Obey the Law}. Princeton University Press.
\bibitem{uduku2002} Uduku, O. (2002). The socio-economic basis of a diaspora community. \textit{African Affairs}, 101(404), 339--355.
\bibitem{un2024} United Nations DESA. (2024). \textit{UN E-Government Survey 2024}. United Nations.
\bibitem{wb2025} World Bank. (2025). \textit{GovTech Maturity Index 2025 update}.
\bibitem{wjp2025} World Justice Project. (2025). \textit{WJP Rule of Law Index 2025}.

\end{thebibliography}
\end{document}
